\documentclass[12pt]{article}
\usepackage[utf8]{inputenc}
\usepackage{geometry}
\usepackage{enumitem}
\usepackage{hyperref}
\usepackage{xcolor}
%\usepackage{fontspec}
% For Malayalam text if available (optional)
% \setmainfont{Noto Serif Malayalam}

\geometry{a4paper, margin=1in}
\hypersetup{
    colorlinks=true,
    linkcolor=blue,
    urlcolor=blue,
}

\title{The Unseen Depths of Kerala: Culture, Manuscripts, and Language}
\author{Your Name}
\date{\today}

\begin{document}

\maketitle

\begin{abstract}
Kerala, often celebrated for its backwaters and high literacy rates, holds layers of cultural and linguistic heritage that remain largely invisible to the mainstream Indian narrative. This document explores three such hidden dimensions: the living traditions of ritualistic craftsmanship, the vast repositories of pre-modern knowledge preserved on palm leaves, and the complex linguistic landscape that extends far beyond the Malayalam language. Drawing on recent ethnographic studies, archival projects, and linguistic surveys, this paper reveals a Kerala defined by its diversity, isolation, and historical depth.
\end{abstract}

\section*{Introduction}
While Kerala is often portrayed as a monolithic cultural entity, the reality is a mosaic of tribal communities, ancient craft guilds, and linguistic minorities whose histories challenge the dominant narrative. This report synthesizes findings from the Endangered Archives Programme, recent linguistic surveys, and ethnographic studies to present a holistic view of Kerala's hidden heritage.

\section{Vibrant Culture: Beyond the Backwaters}
The cultural vibrancy of Kerala is not limited to its renowned festivals like Onam or the performing arts of Kathakali. It is preserved in the dusty workshops of rural villages and the protected groves of its forests.

\subsection{The Moosharis of Kunhimangalam: Crafting Divinity}
In the quiet village of Kunhimangalam in North Kerala, a community known as the Moosharis (bell-metal workers) continues a 900-year-old legacy of lost wax casting. Designated a heritage village in 2018, Kunhimangalam is home to 32 families who create intricate lamps and deities using the ancient \textit{Panchaloha} (five-metal alloy) technique. The process is deeply spiritual: artisans consult Sanskrit scriptures to ensure the correct gestures (\textit{mudras}) of idols and commence work only on auspicious days. This art form, which provides articles for Hindu temples, Christian churches, and Muslim mosques, is a rare example of interfaith craftsmanship rooted in a single community \cite{moosharis}.

\subsection{The Cholanaikkans: A Life Governed by Sunlight}
Hidden in the Nilambur forest lives the Cholanaikkan tribe, one of India's most isolated indigenous communities. Numbering only a few hundred, they are hunter-gatherers who do not rely on calendars or watches. Time for them is measured by the movement of the sun, the shift of seasons, and the behavior of animals. Their existence is a living repository of oral knowledge regarding medicinal plants and forest ecology. Access to their settlements is strictly regulated by authorities to protect their fragile culture from irreversible change, making them a "hidden" world within Kerala \cite{cholanaikkans}.

\section{Manuscripts: Palm-Leaf Repositories of Knowledge}
Before the digital age, Kerala's knowledge was preserved on palm leaves (\textit{tāliyola}) using styluses. These manuscripts, written in scripts like Grantha and Malayalam, are only now being systematically cataloged.

\subsection{The Vadakke Madham Collection}
The Vadakke Madham Brahmaswam in Thrissur houses a collection of 748 palm-leaf manuscripts that represent the intellectual history of the region. These texts were originally accumulated by four monasteries surrounding the Vadakkumnathan Temple. They cover a wide range of subjects, including Vedic philosophy, ritual, and literature. Written primarily in Malayalam and Grantha script, these manuscripts date back centuries, with colophons referencing the Kollam calendar era. The collection is now preserved under strict climate control at the Kerala Manuscript Preservation Centre, a joint effort between the Vedic Research Centre and the École française d’Extrême-Orient \cite{palmleaf}.

\subsection{The Kaṇakkatikāram: Mathematics in Malayalam}
A unique genre of manuscript found in Kerala is the \textit{Kaṇakkatikāram}—elementary mathematical treatises focusing on measurements and practical calculations. Unlike the famous \textit{Kerala School of Astronomy} (which produced infinite series), these texts were pedagogical tools used in medieval and colonial Kerala. They demonstrate how mathematical knowledge was transmitted in the vernacular, mixing word problems with cultural values and professional practices of the time. The manuscripts are notable for their stylistic form and provide insight into the economic realities of pre-modern Kerala \cite{math}.

\section{Local Languages: A Linguistic Mosaic}
While Malayalam is the official language, Kerala is home to several linguistic minorities and scriptorial variants that are often overlooked or marginalized.

\subsection{Mavilan Tulu: A Threatened Language}
A 2025 sociolinguistic survey by SIL International confirmed that Mavilan Tulu, spoken by the Mavilan community (a Scheduled Tribe in Kasargod district), is distinct from both standard Tulu and Malayalam. The study estimates a population of about 45,000 speakers, but notes a disruption in intergenerational transmission. Malayalam is encroaching into informal domains, placing the language at EGIDS level 6b (threatened). However, community organizations like the Akhila Kerala Mavilan Samajam are actively working to preserve the language and culture \cite{mavilan}.

\subsection{Arabi-Malayalam: The Forgotten Script}
A significant part of Kerala's literary heritage exists in Arabi-Malayalam—a scriptorial variant of Malayalam that uses the Arabic script. Popular among the Mappila Muslim community of the Malabar region, this form has historically been overlooked in mainstream literary historiographies. Academic analysis reveals that works in Arabi-Malayalam were often "folklorized" and temporally displaced in historical records, being treated as a lesser form of Malayalam rather than a distinct linguistic tradition. This marginalization has contributed to its current endangered status \cite{arabi}.

\subsection{Minority Languages}
According to census data, while Malayalam is spoken by 97.03\% of the population, minority languages like Tamil (1.49\%), Tulu (0.37\%), Kannada (0.23\%), and Konkani (0.2\%) maintain significant speaker bases. Additionally, isolated communities like the Cholanaikkans speak a distinct Dravidian language influenced by Tamil, Kannada, and Malayalam, but understood nowhere else on Earth \cite{wiki, cholanaikkans}.

\section*{Conclusion}
The culture, manuscripts, and languages of Kerala reveal a state far more complex than its tourist board advertisements suggest. From the bell-metal workers invoking ancient scriptures to the digitization efforts preserving fading scripts, Kerala's heritage is a dynamic interplay of preservation and loss. Acknowledging these hidden elements is crucial for understanding the true fabric of one of India's most unique regions.

\bibliographystyle{plain}
\begin{thebibliography}{9}

\bibitem{moosharis}
Manjulika Pramod, “Kunhimangalam: A Sleepy Town That Keeps A 900-Year Old Legacy Of Bell Metal Crafting Alive,” \textit{Travel and Leisure Asia}, December 30, 2024.

\bibitem{cholanaikkans}
Vivek, “Hidden world of Cholanaikkans, isolated Kerala tribe that lives by sunlight,” \textit{India Today}, December 19, 2025.

\bibitem{palmleaf}
“Palm-leaf manuscripts of the Vadakke Madham Bramaswam, Thrissur,” \textit{Endangered Archives Programme}, British Library (EAP1304).

\bibitem{math}
Roy Wagner and Arun Ashokan, “Content and Context of Kaṇakkatikāram Manuscripts: Pre-Modern Malayalam Elementary Mathematics,” \textit{History of Science in South Asia}, Vol. 12, July 3, 2024.

\bibitem{mavilan}
Maggie Canvin, Nidhin Joseph, and Manoj M. R., “A Sociolinguistic Survey of the Mavilan Tulu Language in Northern Kerala,” \textit{SIL Journal of Language Survey Reports}, 2025-005, 2025.

\bibitem{arabi}
Yoonus Kozhisseri, “Language and Lesser Forms of Language: Arabi-Malayalam Negotiating the Canon with Malayalam Literary Historiographies,” \textit{ICLA Journal}, May 5, 2025.

\bibitem{wiki}
“Ethnic groups in Kerala,” \textit{Wikipedia}, citing Census of India 2011 data.

\end{thebibliography}

\end{document}